\documentclass[11pt]{article}

\usepackage[margin=1in]{geometry}
\usepackage{amsmath,amssymb,amsthm}
\usepackage{graphicx}
\usepackage[hidelinks]{hyperref}

\newtheorem{theorem}{Theorem}
\newtheorem{lemma}[theorem]{Lemma}
\newtheorem{remark}[theorem]{Remark}

\newcommand{\R}{\mathbb R}
\newcommand{\Pbb}{\mathbb P}
\newcommand{\E}{\mathbb E}
\newcommand{\norm}[1]{\left\lVert #1 \right\rVert}
\newcommand{\abs}[1]{\left\lvert #1 \right\rvert}

\title{A Weak-\(\ell_1\) Solution of the \(B_1^n\) Linear Part 1' Step}
\author{arXiv automatic math run \texttt{fa\_banach\_001}}
\date{July 3, 2026}

\begin{document}
\maketitle

\section*{Status}

This packet gives a full solution candidate, likely valid pending expert
review, for the linear Part 1' separation-dimension-reduction challenge in
Mendelson--Milman--Paouris \cite{MMP} in the case \(K=B_1^n\).  The target
body after dimension reduction is a fixed multiple of the weak-\(\ell_1\)
star body.  The result upgrades the earlier ordinary-\(\ell_1\) Cauchy packet:
the ordinary-\(\ell_1\) body incurred an \(O(k\log k)\) separation loss, while
the weak-\(\ell_1\) body gives a universal constant loss.

The theorem below solves the linear Part 1' step, not the entire generalized
dual Sudakov program.  To run the full program with this target family one
would still need the corresponding Part 2 weak Sudakov estimate for
weak-\(\ell_1\) star bodies.

\section*{Source Challenge}

The source paper is Shahar Mendelson, Emanuel Milman, and Grigoris Paouris,
\emph{Generalized Dual Sudakov Minoration via Dimension Reduction - A Program},
arXiv:1610.09287.  The paper formulates a linear version of Part 1:
given \(e^k\) many \(K\)-separated points and a measure \(\mu\) with
\(\mu(K)\ge 1/e\), find a linear map \(T:\R^n\to\R^m\), \(m\le A k\), and a
massive star body \(L\subset\R^m\) such that the images remain
\((1/R)L\)-separated.  In the conclusion, the authors identify Part 1 as the
main remaining challenge and explicitly say that it is significant even for
\(K=B_1^n\).

\begin{center}
\includegraphics[width=\textwidth]{figures/open_problem_crop.png}
\end{center}

\section*{Definitions}

For \(y\in\R^m\), define the weak-\(\ell_1\) quasi-norm
\[
  \norm{y}_{1,\infty}
  =
  \sup_{s>0} s\,\#\{j:\abs{y_j}>s\}.
\]
Equivalently, if \(y_1^*\ge \cdots \ge y_m^*\) is the non-increasing
rearrangement of \((\abs{y_j})\), then
\[
  \norm{y}_{1,\infty}=\max_{1\le r\le m} r y_r^* .
\]
Let
\[
  W_m=\{y\in\R^m:\norm{y}_{1,\infty}\le 1\}.
\]
This is an origin-symmetric compact star body with non-empty interior.

We use the source paper's convention that \(K\)-separated points are points
whose translates by \(K\) are mutually disjoint.  For \(K=B_1^n\), this implies
\(\norm{x_i-x_j}_1\ge 2\) up to the immaterial boundary convention; the proof
below only uses the weaker lower bound \(\norm{x_i-x_j}_1\ge 1\).

\section*{Proof Intuition}

The Cauchy law is exactly \(1\)-stable: a row of a Cauchy matrix applied to
\(x\) has scale \(\norm{x}_1\).  In the earlier ordinary-\(\ell_1\) attempt,
the vector of \(m\) Cauchy variables had typical \(\ell_1\)-size
\(m\log m\), which produced the unwanted \(O(k\log k)\) loss.

The right scale for Cauchy tails is weak \(\ell_1\).  If
\(G=(G_1,\ldots,G_m)\) has independent standard Cauchy coordinates, then
\(\norm{G}_{1,\infty}=O(m)\) with fixed probability.  After scaling the random
map by \(1/m\), the image of every point of \(B_1^n\) falls in a fixed multiple
of \(W_m\) with fixed probability, yielding massiveness by averaging.  At the
same time, the event \(\norm{G}_{1,\infty}\le \eta m\) has exponentially small
probability for small \(\eta\), because it would force far fewer than the
typical number of coordinates to exceed \(1\).  This exponential small-ball
bound is strong enough to union-bound over \(e^k\) point pairs in dimension
\(m\simeq k\).

\section*{The Result}

\begin{theorem}[Weak-\(\ell_1\) Part 1' for \(B_1^n\)]
\label{thm:main}
There are universal constants \(A_0,A_1,R_0,q_0,K_0>0\) with
\(q_0\ge 1\) such that the following holds.  Let \(k\ge K_0\), let
\(M\le e^k\), and let \(x_1,\ldots,x_M\in\R^n\) satisfy
\[
  \norm{x_i-x_j}_1\ge 1 \qquad (i\ne j).
\]
Let \(\mu\) be a Borel probability measure on \(\R^n\) with
\(\mu(B_1^n)\ge 1/e\).  Then there are \(m\le A_1 k\), a linear map
\(T:\R^n\to\R^m\), and
\[
  L=A_0 W_m
\]
such that
\[
  T x_i - T x_j \notin \frac{1}{R_0}L
  \qquad (i\ne j)
\]
and
\[
  T_*\mu(L)\ge e^{-q_0}.
\]
In particular, this verifies the linear Part 1' separation-dimension-reduction
step for \(K=B_1^n\), with universal constants, using the reduced target
family \(L=A_0W_m\).
\end{theorem}

\begin{lemma}[Weak-\(\ell_1\) estimates for Cauchy vectors]
\label{lem:cauchy-weak}
Let \(G_1,\ldots,G_m\) be independent standard Cauchy random variables and set
\(G=(G_1,\ldots,G_m)\).  There are universal constants
\(\eta,A,c>0\) such that, for every \(m\ge 1\),
\[
  \Pbb\{\norm{G}_{1,\infty}\le \eta m\}\le e^{-cm}
\]
and
\[
  \Pbb\{\norm{G}_{1,\infty}\le A m\}\ge \frac34 .
\]
\end{lemma}

\begin{proof}
For the lower-tail estimate, use
\(\Pbb\{\abs{G_1}>1\}=1/2\).  If \(\norm{G}_{1,\infty}\le \eta m\), then
\[
  \#\{j:\abs{G_j}>1\}\le \eta m.
\]
Taking \(\eta>0\) sufficiently small, the Chernoff bound for a binomial
\(\mathrm{Bin}(m,1/2)\) variable gives
\[
  \Pbb\{\norm{G}_{1,\infty}\le \eta m\}\le e^{-cm}.
\]

For the upper-tail estimate, write
\[
  N(u)=\#\{j:\abs{G_j}>u\}.
\]
There is a universal \(C_0\) such that
\(\Pbb\{\abs{G_1}>u\}\le C_0/u\) for all \(u>0\).  If
\(\norm{G}_{1,\infty}>A m\), then for some \(u>A\) one has
\(uN(u)>A m\).  Choose \(r\ge 0\) with
\(2^r A<u\le 2^{r+1}A\).  Then
\[
  N(2^r A)> \frac{m}{2^{r+1}}.
\]
For those \(r\) with \(m/2^{r+1}\ge 1\), the binomial tail bound gives
\[
\Pbb\left\{N(2^r A)> \frac{m}{2^{r+1}}\right\}
\le
\left(\frac{2eC_0}{A}\right)^{m/2^{r+1}},
\]
provided \(A\) is larger than a fixed multiple of \(C_0\).  Summing over these
dyadic \(r\)'s yields at most
\[
  \sum_{s=1,2,4,\ldots}\left(\frac{2eC_0}{A}\right)^s.
\]
The remaining dyadic levels have \(2^rA\gtrsim Am\), and their total
contribution is bounded by
\[
  \Pbb\{\max_j \abs{G_j}\gtrsim A m\}\le \frac{C_1}{A}.
\]
Choosing \(A\) sufficiently large makes the total probability at most \(1/4\).
\end{proof}

\begin{proof}[Proof of Theorem~\ref{thm:main}]
Let \(m=\lceil C_2 k\rceil\), where \(C_2\) will be chosen as a sufficiently
large universal constant.  Let \((g_{jr})_{1\le j\le m,1\le r\le n}\) be
independent standard Cauchy variables and define
\[
  (Tx)_j=\frac1m\sum_{r=1}^n g_{jr}x_r,
  \qquad j=1,\ldots,m.
\]
The standard \(1\)-stability identity for Cauchy variables gives, for each
fixed \(x\in\R^n\),
\[
  Tx \stackrel{d}{=} \frac{\norm{x}_1}{m}(G_1,\ldots,G_m),
\]
with independent standard Cauchy \(G_j\)'s.

Let \(h=x_i-x_j\).  Since \(\norm{h}_1\ge 1\), Lemma
\ref{lem:cauchy-weak} implies
\[
  \Pbb\{\norm{Th}_{1,\infty}\le \eta\}
  \le
  \Pbb\{\norm{G}_{1,\infty}\le \eta m\}
  \le e^{-cm}.
\]
Therefore
\[
  \Pbb\{\exists i<j:\norm{T(x_i-x_j)}_{1,\infty}\le \eta\}
  \le M^2 e^{-cm}\le \exp(2k-cm).
\]
Choose \(C_2\) so large that this probability is at most \(e^{-3}\).

Now set \(L=A W_m\), where \(A\) is the constant in Lemma
\ref{lem:cauchy-weak}.  For every \(x\in B_1^n\), the same stability identity
and the upper-tail estimate give
\[
  \Pbb\{Tx\in L\}\ge \frac34 .
\]
Consequently,
\[
  \E\bigl[T_*\mu(L)\bigr]
  =
  \int \Pbb\{Tx\in L\}\,d\mu(x)
  \ge
  \frac34\,\mu(B_1^n)
  \ge
  \frac{3}{4e}.
\]
Since \(0\le T_*\mu(L)\le 1\), it follows that
\[
  \Pbb\left\{T_*\mu(L)\ge \frac{3}{8e}\right\}\ge \frac{3}{8e}.
\]
After increasing \(C_2\) if needed, the separation failure probability is
smaller than \(3/(8e)\).  Hence there exists a realization of \(T\) for which
both events hold:
\[
  \norm{T(x_i-x_j)}_{1,\infty}>\eta \quad (i\ne j),
  \qquad
  T_*\mu(AW_m)\ge \frac{3}{8e}.
\]
Put \(A_0=A\), \(R_0=A/\eta\), and
\[
  q_0=\log\left(\frac{8e}{3}\right).
\]
Then \(T_*\mu(L)\ge e^{-q_0}\), and
\[
  \eta W_m=\frac{\eta}{A}L=\frac1{R_0}L,
\]
so the images are \((1/R_0)L\)-separated.  Finally, for
\(K_0\ge 2q_0\) we have \(q_0\le k/2\), exactly as required in Part 1'.
\end{proof}

\section*{Verification and Novelty Check}

The proof uses only the elementary \(1\)-stability of Cauchy variables,
Chernoff bounds for binomial variables, and a dyadic weak-\(\ell_1\) tail
estimate.  No computational verification is required.

Local run indexes were searched for arXiv:1610.09287, ``Cauchy'',
``weak l1'', ``weak-\(\ell_1\)'', ``Part 1'', and
``one-sided Johnson--Lindenstrauss''.  They contained the earlier ordinary
\(\ell_1\) partial packet but no weak-\(\ell_1\) full Part 1' packet.
External web searches on July 3, 2026 for combinations of
``Cauchy weak l1 dimension reduction'', ``one-sided Johnson-Lindenstrauss
Cauchy weak l1'', and ``B1 weak l1 Sudakov dimension reduction'' did not
surface a direct match.  The arXiv abstract for \cite{MMP} records that the
paper establishes the separation step for ellipsoids and, with logarithmic
losses, for cubes; it does not state the weak-\(\ell_1\) \(B_1^n\) result
proved here.

\section*{Limitations and Human Review Recommendation}

This packet should be reviewed as a full solution of the linear Part 1' step
for \(K=B_1^n\).  It should not be cited as a full proof of the generalized
dual Sudakov conjecture, because Part 2 has not been proved for the
weak-\(\ell_1\) target family \(A W_m\).

The most important checks are:
\begin{itemize}
\item the dyadic proof that \(\norm{G}_{1,\infty}=O(m)\) with constant
probability;
\item the lower-tail estimate \(\Pbb\{\norm{G}_{1,\infty}\le\eta m\}\le
e^{-cm}\);
\item the translation from \(\norm{T(x_i-x_j)}_{1,\infty}>\eta\) to
\((1/R_0)L\)-separation.
\end{itemize}

\begin{thebibliography}{9}

\bibitem{MMP}
S. Mendelson, E. Milman, and G. Paouris,
\emph{Generalized Dual Sudakov Minoration via Dimension Reduction - A Program},
arXiv:1610.09287, 2016; revised 2018.

\end{thebibliography}

\end{document}
